\section*{Exercises about Positive Operators}

\exercise{1}
\begin{xrcs}
  Suppose $T \in \linmap(V)$. Prove that if both $T$ and $-T$ are positive operators, then $T = 0$.

\begin{xprf}
  Since $T$ is positive, $\langle Tv, v \rangle \geq 0$ for all $v \in V$.
  Since $-T$ is positive, $\langle -Tv, v \rangle = -\langle Tv, v \rangle \geq 0$, which implies $\langle Tv, v \rangle \leq 0$ for all $v \in V$.
  Therefore, $\langle Tv, v \rangle = 0$ for all $v \in V$.
  Since $T$ is positive, it is self-adjoint.
  By Theorem \ref{thm: T self-adjoint and <Tv,v>=0 for all v <=> T=0}, $\langle Tv, v \rangle = 0$ for all $v \in V$ implies $T = 0$.
\end{xprf}
\end{xrcs}

\exercise{6}
\begin{xrcs}
  Prove that the sum of any two positive operators on $V$ is a positive operator.

\begin{xprf}
  Let $S, T \in \linmap(V)$ be positive operators.
  By definition of a positive operator, $S$ and $T$ are self-adjoint, and $\langle Sv, v \rangle \geq 0$ and $\langle Tv, v \rangle \geq 0$ for all $v \in V$.
  First, $(S+T)^* = S^* + T^* = S + T$, so $S+T$ is self-adjoint.
  Second, for every $v \in V$:
  \[
    \langle (S+T)v, v \rangle = \langle Sv + Tv, v \rangle = \langle Sv, v \rangle + \langle Tv, v \rangle \geq 0 + 0 = 0.
  \]
  Therefore, $S+T$ is a positive operator.
\end{xprf}
\end{xrcs}

\exercise{12}
\begin{xrcs}
  Suppose $T \in \linmap(V)$ is a positive operator. Prove that $T^k$ is a positive operator for every positive integer $k$.

\begin{xprf}
  Since $T$ is positive, it is self-adjoint.
  By the Spectral Theorem (Theorem \ref{thm: real spectral theorem} or \ref{thm: complex spectral theorem}), there exists an orthonormal basis $e_1, \dots, e_n$ of $V$ consisting of eigenvectors of $T$.
  Let $\lambda_1, \dots, \lambda_n$ be the corresponding eigenvalues.
  By Theorem \ref{thm: characterizations of positive operators}(b), $\lambda_j \geq 0$ for each $j \in \{1, \dots, n\}$.
  Then $T^k e_j = \lambda_j^k e_j$ for each $j$.
  Since each $\lambda_j \geq 0$, we have $\lambda_j^k \geq 0$.
  Thus with respect to the orthonormal basis $e_1, \dots, e_n$, $\mmatrix(T^k)$ is a diagonal matrix whose diagonal entries are all non-negative numbers $\lambda_j^k \geq 0$.
  By Theorem \ref{thm: characterizations of positive operators}(c), $T^k$ is a positive operator.
\end{xprf}
\end{xrcs}

\exercise{16}
\begin{xrcs}
  Suppose $T$ is a positive operator on $V$. Prove that
  \begin{equation}
    \mynull \sqrt{T} = \mynull T \quad \text{and} \quad \myrange \sqrt{T} = \myrange T.
  \end{equation}

\begin{xprf}
  \StepOne If $v \in \mynull \sqrt{T}$, then $\sqrt{T}v = 0 \implies Tv = \sqrt{T}(\sqrt{T}v) = 0$, so $\mynull \sqrt{T} \subseteq \mynull T$.
  Conversely, if $v \in \mynull T$, then $Tv = 0$, so $\langle Tv, v \rangle = 0$.
  Since $T$ is positive, $0 = \langle Tv, v \rangle = \langle \sqrt{T}\sqrt{T}v, v \rangle = \langle \sqrt{T}v, \sqrt{T}v \rangle = \|\sqrt{T}v\|^2 \implies \sqrt{T}v = 0$.
  Thus $\mynull T \subseteq \mynull \sqrt{T}$, so $\mynull \sqrt{T} = \mynull T$.

  \StepTwo Since both $T$ and $\sqrt{T}$ are self-adjoint, by Theorem \ref{thm: null space and range of T^*}:
  \[
    \myrange \sqrt{T} = \left(\mynull \left(\sqrt{T}\right)^*\right)^\perp = (\mynull \sqrt{T})^\perp = (\mynull T)^\perp = (\mynull T^*)^\perp = \myrange T.
  \]
\end{xprf}
\end{xrcs}
